\chapter{Counting technique}

\begin{orient}
\textbf{What this chapter is for.} Counting is a translation problem. There is a small stock of sets whose size everyone knows, and every counting technique is a way of exhibiting the set in front of you as one of them in disguise. A bijection says the two sets are the same set relabelled. Double counting says one set has two descriptions and therefore one size with two formulas. Inclusion and exclusion says an overcount can be repaired systematically. Complementary counting says the easier set to count is the one you were not asked about. Recursion says the set decomposes into smaller copies of itself. This chapter treats each as a recognition problem, and every count in it has been checked against brute-force enumeration for small parameters.
\end{orient}

\section{Making the set into a set you know}

\tech{Bijective counting}{medium}
\cue Two sets that ought to have the same size, or a set whose elements are naturally encoded as strings, paths, subsets, or arrangements. Also any answer that comes out as a familiar binomial coefficient, which is a hint that a bijection was available.
\method Exhibit an explicit map, verify it is injective by describing the inverse, and count the target instead. A bijection is the strongest kind of counting argument because it survives generalisation: the same map usually answers the next four questions in the family.

\begin{worked}
Count the lattice paths from $(0,0)$ to $(m,n)$ using unit steps right and up. Encode a path as the string of its steps, a word of length $m+n$ in which $m$ letters are R. The map from paths to such words is a bijection, and choosing which positions hold the R's gives $\binom{m+n}{m}$. For $m=3$, $n=4$ this is $35$, confirmed by enumerating all $2^{7}$ words and keeping those with three R's.
\end{worked}

\begin{trap}
Describing a map and never checking injectivity. The usual failure is a map that forgets an ordering, so that several objects share an image; the repair is either to divide by the size of the fibre, which requires the fibres to be equal in size, or to enrich the target.
\end{trap}

\begin{seealsob}Stars and bars; reflection and the ballot problem; double counting.\end{seealsob}

\tech{Double counting}{medium}
\cue An identity to prove, especially one with binomial coefficients on both sides; or a quantity that can be totalled by rows and by columns; or a problem about incidences between two kinds of object.
\method Find a single set and count it two ways. Equating the two counts proves the identity, and the proof carries an explanation that an algebraic manipulation does not. The standard frame is a table: sum the entries by row and then by column.

\begin{worked}
Prove $\sum_{k=0}^{n}k\binom{n}{k}=n2^{n-1}$. Count the pairs $(S,x)$ where $S\subseteq\{1,\ldots,n\}$ and $x\in S$. By choosing $S$ first, the count is $\sum_{k}k\binom{n}{k}$. By choosing $x$ first, there are $n$ choices, and then any of the $2^{n-1}$ subsets of the remaining elements may be adjoined, giving $n2^{n-1}$. Checked by direct summation for $n$ up to $12$.
\end{worked}

\begin{trap}
Counting two different sets. Write down the set once, in words, and then count that same sentence twice; if the two descriptions do not both describe it, the identity you prove is not the one you wanted.
\end{trap}

\begin{seealsob}Bijective counting; the handshake identity; Burnside and symmetry counting.\end{seealsob}

\section{Repairing an overcount}

\tech{Inclusion and exclusion}{medium}
\cue A count of objects avoiding several forbidden properties, where each property alone is easy to count but they overlap. The words ``at least one'', ``none of'', and ``exactly $k$ of'' are all cues.
\method Add the sizes of the sets, subtract the pairwise intersections, add the triples, and so on:
\[
\abs{A_{1}\cup\cdots\cup A_{n}}=\sum_{\emptyset\neq S\subseteq[n]}(-1)^{\abs{S}+1}\Bigl|\bigcap_{i\in S}A_{i}\Bigr|.
\]
The sign discipline is the whole difficulty, and the practical form is: count the complement, so that the alternating sum runs over ``how many of the bad properties hold''.

\begin{worked}
Derangements. Let $A_{i}$ be the permutations of $[n]$ fixing $i$. Then $\abs{\bigcap_{i\in S}A_{i}}=(n-\abs{S})!$, so
\[
D_{n}=n!\sum_{k=0}^{n}\frac{(-1)^{k}}{k!},
\]
giving $D_{4}=9$ and $D_{5}=44$, both confirmed by exhaustive enumeration. The ratio $D_{n}/n!$ tends to $1/\eu$, which is why the probability that a random permutation has no fixed point is close to $0.3679$ for every $n$ beyond about six.
\end{worked}

\begin{trap}
Applying the formula when the intersections do not have a size depending only on $\abs{S}$. In the derangement case they do, which is what collapses $2^{n}$ terms into $n+1$; in general you must sum over the subsets themselves.
\end{trap}

\begin{seealsob}Complementary counting; recursion and generating functions; Burnside.\end{seealsob}

\tech{Complementary counting}{easy}
\cue The words ``at least one'' in the question. Also any count where the forbidden configurations are few and the permitted ones are many.
\method Count everything and subtract what fails. The test for whether this is the right move is simple: compare the number of cases in the direct count with the number in the complement, and take the smaller.

\begin{worked}
How many five-card hands from a standard deck contain at least one ace? Directly this is a four-case sum. By complement it is $\binom{52}{5}-\binom{48}{5}=2598960-1712304=886656$.
\end{worked}

\begin{seealsob}Inclusion and exclusion; bijective counting.\end{seealsob}

\section{Distributions and recursions}

\tech{Stars and bars, compositions, and the hockey stick}{easy}
\cue Distributing identical objects into distinct boxes, or counting non-negative integer solutions of $x_{1}+\cdots+x_{k}=n$, or a sum of binomial coefficients along a diagonal.
\method Non-negative solutions number $\binom{n+k-1}{k-1}$; positive solutions number $\binom{n-1}{k-1}$. The proof is the bijection with arrangements of $n$ stars and $k-1$ bars. The hockey-stick identity $\sum_{i=r}^{n}\binom{i}{r}=\binom{n+1}{r+1}$ is the same idea read along a diagonal of Pascal's triangle, and it is the standard way to collapse a sum of binomials.

\begin{worked}
Non-negative integer solutions of $x_{1}+x_{2}+x_{3}=7$: $\binom{9}{2}=36$, confirmed by enumeration. With the constraint $x_{1}\geq2$, substitute $x_{1}'=x_{1}-2$ to get $x_{1}'+x_{2}+x_{3}=5$ and $\binom{7}{2}=21$. Upper bounds are handled by inclusion and exclusion on the violated constraints rather than by substitution.
\end{worked}

\begin{trap}
Using the identical-objects formula for distinguishable objects. Seven distinct balls into three distinct boxes is $3^{7}$, not $\binom{9}{2}$, and the two are confused whenever the problem statement is read quickly.
\end{trap}

\begin{seealsob}Bijective counting; inclusion and exclusion.\end{seealsob}

\tech{Recursion, state, and generating functions}{hard}
\cue A count parameterised by $n$ where the objects of size $n$ decompose into objects of smaller size, especially with a forbidden local pattern (``no two adjacent'', ``no three in a row'').
\method Define the count with enough state to make the recursion close. The commonest error is too little state, and the fix is to split by what the last element was. Then either solve the recurrence by the characteristic equation, or encode it as a generating function $F(x)=\sum a_{n}x^{n}$, translate the recursion into a functional equation, and extract coefficients.

\begin{worked}
Binary strings of length $n$ with no two consecutive ones. Let $a_{n}$ be the count. A valid string ends in $0$, preceded by any valid string of length $n-1$, or ends in $01$, preceded by any valid string of length $n-2$. So $a_{n}=a_{n-1}+a_{n-2}$ with $a_{1}=2$, $a_{2}=3$, giving the Fibonacci numbers shifted: $a_{n}=F_{n+2}$. For $n=6$ this predicts $21$, confirmed by enumerating all $64$ strings.
\end{worked}

\begin{worked}[The same count by generating function]
Every valid string is a sequence of blocks, each block being $0$ or $10$, with an optional leading $1$. The generating function is therefore $\frac{1+x}{1-x-x^{2}}$, whose denominator reproduces the Fibonacci recurrence directly. Reading the structure off the object and reading the recurrence off the structure are the same step.
\end{worked}

\begin{trap}
A recursion whose base cases do not match the convention used in the step. Compute $a_{1}$ and $a_{2}$ by hand, then check $a_{3}$ against the recursion before trusting anything.
\end{trap}

\begin{seealsob}Induction; reflection and the ballot problem; series evaluation by generating function.\end{seealsob}

\tech{Reflection, the ballot problem, and Catalan structure}{hard}
\cue Paths constrained to stay on one side of a line, balanced bracket sequences, non-crossing configurations, binary trees, triangulations of a polygon. The number $\binom{2n}{n}/(n+1)$ appearing anywhere is the signature.
\method The reflection principle counts bad paths by reflecting the portion before the first violation, which biject bad paths to unconstrained paths with shifted endpoints. Subtracting gives
\[
C_{n}=\frac{1}{n+1}\binom{2n}{n},
\]
the Catalan numbers $1,1,2,5,14,42,132,\ldots$ Recognising a Catalan family is worth more than rederiving the formula, because the same numbers answer at least a dozen superficially unrelated questions.

\begin{worked}
Balanced strings of $n$ pairs of brackets. Encode an open bracket as an up step and a close as a down step. Balanced means the path ends at zero and never goes below it. Total paths $\binom{2n}{n}$; bad paths reflect to paths ending at $-2$, of which there are $\binom{2n}{n+1}$. The difference is $\binom{2n}{n}-\binom{2n}{n+1}=C_{n}$. For $n=4$ this gives $14$, confirmed by enumerating all $70$ balanced-length strings and testing each prefix.
\end{worked}

\begin{seealsob}Bijective counting; recursion and generating functions.\end{seealsob}

\tech{Burnside and symmetry counting}{hard}
\cue Counting configurations up to a symmetry: colourings of a necklace up to rotation, of a cube's faces up to rotation, of a board up to reflection. The cue is the phrase ``considered the same if''.
\method The number of orbits is the average number of fixed points of the group elements:
\[
\abs{X/G}=\frac{1}{\abs{G}}\sum_{g\in G}\abs{\mathrm{Fix}(g)}.
\]
The work is entirely in computing $\abs{\mathrm{Fix}(g)}$ for each symmetry, which for a rotation by $d$ positions out of $n$ with $c$ colours is $c^{\gcd(d,n)}$.

\begin{worked}
Necklaces of six beads in three colours, up to rotation. The six rotations have $\gcd(d,6)=6,1,2,3,2,1$ for $d=0,\ldots,5$, so the fixed-point counts are $3^{6},3,3^{2},3^{3},3^{2},3$, totalling $729+3+9+27+9+3=780$, and $780/6=130$. Confirmed by generating all $729$ colourings and counting rotation classes.
\end{worked}

\begin{trap}
Dividing the total by the group size without averaging. That shortcut is correct only when no configuration has a nontrivial symmetry, which is exactly the case the problem is usually designed to break.
\end{trap}

\begin{seealsob}Double counting; bijective counting; symmetry reduction.\end{seealsob}

\section{Recognition matrix}

\begin{recog}{@{}p{0.31\textwidth}p{0.35\textwidth}p{0.28\textwidth}@{}}
\rowcolor{mist}\mh{If you see} & \mh{Reach for} & \mh{If that fails} \\
\midrule
\endhead
Two sets that should match & An explicit bijection & Count both directly \\
A binomial identity to prove & Double counting one set & Induction on $n$ \\
``At least one'' & Complementary counting & Inclusion and exclusion \\
Several overlapping forbidden traits & Inclusion and exclusion & Group by how many hold \\
No fixed points, no repeats & Derangements, $D_n=n!\sum(-1)^k/k!$ & Recursion $D_n=(n-1)(D_{n-1}+D_{n-2})$ \\
Identical items, distinct boxes & Stars and bars & Substitute away lower bounds \\
Distinct items, distinct boxes & $k^{n}$, or surjections by IE & Exponential generating function \\
A diagonal sum of binomials & Hockey stick & Absorb and telescope \\
``No two adjacent'' & Recursion split by the last element & Transfer matrix \\
Paths staying above a line & Reflection; expect Catalan & Ballot formula \\
$\binom{2n}{n}/(n+1)$ appears & You are in a Catalan family & Find the bracket encoding \\
``Considered the same if rotated'' & Burnside averaging & Fix a canonical representative \\
\end{recog}
